\chapter[Falsifiers]{Assumptions, Predictions, and Falsifiers}
\label{ch:falsifiers}

\paragraph{Learning outcomes.}
After this chapter the reader should be able to:
\begin{enumerate}
  \item restate the derivation chain of Chapter~1 with the theorem
        numbers now attached;
  \item separate assumptions that kill NAMCG if they fail from
        coefficients that can be adjusted;
  \item name observations that would refute the framework rather
        than merely bound \(\eta\) or \(r_0\);
  \item refuse to treat a successful laboratory plot as evidence
        for Theorems~\ref{thm:mT}--\ref{thm:Leff};
  \item enter Chapter~17 knowing which open calculations are
        optional refinements and which are missing load-bearing
        proofs.
\end{enumerate}

A framework that cannot say what would kill it is not a framework.
This chapter is the scorecard promised in Section~\ref{sec:ch1-checklist}.

\section{The chain, with theorem numbers}
\label{sec:ch16-chain}

The argument of the book is the following sequence, and no shorter
sequence.

\begin{enumerate}
  \item A matrix order parameter \(\Phi^I{}_A\) with potential
        \eqref{eq:V} admits a frame-locking vacuum
        (Assumptions~\ref{ass:frame-lock}, \ref{ass:dictionary}).
  \item A composite vierbein and metric exist in the condensed
        phase (Repair~II,~\eqref{eq:vev-is-frame};
        Definition~\ref{def:composite-metric},
        Assumption~\ref{ass:nondeg}).
  \item Fluctuations split into \(\psi\), \(A_\mu\), \(\omega_{AB}\),
        and \(T_{\mu\nu}\) (Chapter~5).
  \item Entanglement equilibrium on emergent diamonds is a property
        of the condensate vacuum (Assumption~\ref{ass:EE}).
  \item A Fierz--Pauli tensor block is incompatible with that
        equilibrium unless \(m_T=0\) (Theorem~\ref{thm:mT}).
  \item The Euclidean curvature coefficient obeys the conditional,
        truncated heat-kernel identity of Theorem~\ref{thm:Geff};
        physical matching requires specified operators, weights,
        thresholds, continuation, and scheme (Problem~F).
  \item Four-derivative metric operators are EFT vertices, not
        extra poles (Assumption~\ref{ass:EFT-a2}).
  \item Local Lorentz invariance is a gauge symmetry and
        \(\omega_{AB}\) is eaten (Theorem~\ref{thm:locking}).
  \item The conformal mode is heavy (Assumption~\ref{ass:gap-psi}).
  \item The total stress, including \(T_{\mu\nu}^{\mathrm{ent}}\),
        is covariantly conserved (Theorem~\ref{thm:conservation}).
  \item \(\Lambda_{\mathrm{eff}}\) is set by boundary entanglement
        after an ultraviolet subtraction (Theorem~\ref{thm:Leff}).
\end{enumerate}

Items 5, 6, 8, 10, and 11 are theorems \emph{inside} the list of
assumptions that precede them. If an earlier assumption fails, the
later theorem is idle. Phenomenology in Chapters~13--15 sits after
the whole list. It cannot substitute for any line.

\section{Inventory of named assumptions}
\label{sec:ch16-inventory}

The named hypotheses of Parts~II--V are collected here. ``Load-
bearing'' means that failure returns the infrared theory to analog
gravity, scalar-tensor gravity, or an inconsistent extra force.
``Coefficient'' means that failure changes a number or a regime of
validity without destroying the graviton.

\begin{center}
\small
\begin{tabular}{@{}llp{0.38\linewidth}@{}}
\toprule
Tag & Kind & Content \\
\midrule
\ref{ass:aux-metric} & setup & \(g^{(0)}\) is a scaffold \\
\ref{ass:frame-lock} & load-bearing & vacuum
  \(\langle\Phi^I{}_A\rangle=v\delta^I{}_A\) exists \\
\ref{ass:dictionary} & load-bearing & commutator / covariant-derivative
  dictionary \\
\ref{ass:nondeg} & load-bearing & condensed frame is non-degenerate \\
\ref{ass:lorentzian} & load-bearing & signature is Lorentzian \\
\ref{ass:locked-block} & coefficient & extra internal scalars are heavy \\
\ref{ass:EE} & load-bearing & microscopic entanglement equilibrium \\
\ref{ass:gap-psi} & load-bearing & conformal mode is gapped \\
\ref{ass:laplace} & setup & infrared operator is Laplace type \\
\ref{ass:EFT-a2} & load-bearing & \(C^2,R^2\) are vertices, not poles \\
\ref{ass:ECG} & model & cubic truncation in the static spherical
  sector \\
\ref{ass:local-lorentz} & load-bearing & local \(\mathrm{SO}(3,1)\)
  is gauged \\
\ref{ass:Sent-density} & load-bearing & local boundary entropy
  density exists \\
\ref{ass:bdy-EE} & load-bearing & boundary stationarity sets
  \(\Lambda_{\mathrm{eff}}\) \\
\ref{ass:track} & coefficient & static extra density tracks baryons;
  disks use \(\nu_{\sqrt{}}\) of Appendix~\ref{app:empirical} \\
\ref{ass:efe} & coefficient & external field saturates \(\nu\) at \(\nu(g_{\mathrm{ext}})\),
  not \(\nu(g_{\mathrm{N}})\) \\
\ref{ass:drag} & coefficient & working laboratory law
  \(\mathbf{F}_{\mathrm{diss}}\propto -\mathbf{v}/d^4\) \\
\bottomrule
\end{tabular}
\end{center}

Two comments on the table.

First, ``coefficient'' is not a synonym for ``unimportant.''
Assumption~\ref{ass:track} is how galaxies enter the book. If it
fails, galactic claims fail. The graviton claims need not.

Second, Assumption~\ref{ass:dictionary} is the most important
unproved microscopic bridge. If the infrared dictionary from
Presentation~A to Presentation~B cannot be constructed, the heat
kernel of Chapter~8 is a calculation in a different theory than
the matrix action of Chapter~3.

\section{What would kill the framework}
\label{sec:ch16-kill}

The following outcomes are incompatible with NAMCG as written,
not merely with a choice of \(\eta\).

\paragraph{A light extra tensor or a ghost below the cutoff.}
If precision gravity, pulsar timing, or gravitational-wave
propagation requires a second massless or light spin-2 field, or
a ghost pole below \(\LUV\), Theorem~\ref{thm:mT} plus
Assumption~\ref{ass:EFT-a2} has failed.

\paragraph{A light conformal partner.}
If Solar-System or binary-pulsar tests require a Brans--Dicke
parameter that cannot be produced by a heavy \(\psi\), 
Assumption~\ref{ass:gap-psi} has failed and the infrared theory
is scalar-tensor gravity. That is the ceiling of Chapter~1,
reopened.

\paragraph{Non-conservation of the extra stress.}
If an extra acceleration is demanded by data and cannot be written
as \(\nabla^\mu T_{\mu\nu}^{\mathrm{ent}}\) with
Theorem~\ref{thm:conservation} intact---a drag into nowhere, a
preferred-frame force that cannot be improved to a tensor---then
the geometric reading has failed. Analog dissipation remains
possible. Geometrogenesis of the extra channel has not occurred.

\paragraph{Failure of Lorentz locking.}
If the infrared spectrum contains six extra Goldstone polarizations
from \(\omega_{AB}\), Theorem~\ref{thm:locking} has failed. No
later chapter can hide them.

\paragraph{A cutoff that cannot serve Newton and Lorentz together.}
If the only way to obtain \(\Geff=G_{\mathrm{N}}\) forces
\(\LUV\) into a regime where Lorentz-violating \(\Theta^{\mu\nu}\)
corrections are already excluded by current bounds, the
coefficient map and the kinematic algebra cannot coexist.

\paragraph{An induced Einstein--Hilbert term of the wrong sign.}
If a complete heat-kernel evaluation on the locked bundle produces
a repulsive Einstein term that cannot be flipped by fermions or
by the Lorentzian continuation, the infrared force is not gravity.

None of these killers is a rotation curve. None of them is a
pretty laboratory inspiral.

\section{What would only constrain a coefficient}
\label{sec:ch16-constrain}

The following outcomes bound parameters or restrict regimes.
They do not, by themselves, refute Theorems~\ref{thm:mT}--\ref{thm:Leff}.

\begin{itemize}
  \item SPARC ranking of a global \(\nu(g_{\mathrm{N}})\) against
        simple MOND and the published RAR
        (Appendix~\ref{app:empirical}). Frozen
        \(\nu_{\sqrt{}} \) at \(g_\dagger\) is equivalent to the
        RAR and about \(0.2\%\) behind simple MOND in \(\chi^2\);
        baryons-only is not competitive. That is a passed
        coefficient check on disks, not a theorem. A later
        sample that rejects every high-\(g\)-Newtonian,
        low-\(g\)-square-root boost would restrict
        Assumption~\ref{ass:track} to a smaller class of
        systems, or kill the galactic claim only.
  \item Wide binaries in the solar neighbourhood.
        Isolated \(\nu_{\sqrt{}}(g_{\mathrm{N}})\) predicts
        \(\gamma\approx 4\) at \(g_{\mathrm{N}}=10^{-11}\,\mathrm{m\,s^{-2}}\).
        Assumption~\ref{ass:efe} predicts
        \(\gamma\to\nu_{\sqrt{}}(g_{\mathrm{gal}})\approx 1.1\).
        A quality-controlled Gaia sample consistent with Newton
        constrains the isolated reading and is compatible with
        Assumption~\ref{ass:efe}.
        A surviving \(\gamma\approx 1.4\) is not explained by that
        saturation.
        Neither outcome touches Theorems~7.3--12.4.
  \item Cluster lensing or the bullet-cluster mass offset.
        Same status: the extra tensor may fail as a complete
        replacement for particulate dark matter while the graviton
        derivation stands.
  \item Absence of an extra drag in LIGO-band vacuum inspirals.
        Expected if \(\eta\) is negligible at those separations;
        Section~\ref{sec:ch14-quad} already said so.
  \item Absence of a visible core scale \(r_0\) in EHT images or
        ringdowns. Expected if \(r_0\) is microscopic; bounds
        \(\beta_{\mathrm{ECG}}\).
  \item A measured \(\Lambda_{\mathrm{eff}}\) that does not match
        a particular guess for \(S_{\mathrm{ent}}^{\mathrm{(IR)}}\).
        Theorem~\ref{thm:Leff} changes the origin of the scale. It
        does not compute the number.
  \item A different static mutual-information tail does not, by
        itself, determine the drag exponent. The response exponent
        must be calculated from the dynamical force response. The
        \(d^{-4}\) law remains a falsifiable laboratory assumption
        (Assumption~\ref{ass:drag}).
\end{itemize}

\section{Observational checklist}
\label{sec:ch16-obs}

A short list of current and near-term measurements, with the
question each actually asks of the book.

\paragraph{Solar System and pulsar timing.}
Do they tolerate a massive \(\psi\) and a massive torsion at
scale \(v\), and a \(\mu_{\mathrm{ent}}\) that is negligible at
high acceleration? If no, Assumption~\ref{ass:gap-psi} or
\ref{ass:track} dies. If yes, nothing about diamonds has been
tested.

\paragraph{Gravitational waves.}
Is the propagating mode two TT polarizations at the speed of
light, with no extra light tensor? That is a test of
Theorem~\ref{thm:mT} plus Assumption~\ref{ass:EFT-a2}, and it is
already passed at present precision. A measured extra
polarization would be a killer. A measured inspiral-time
deviation at LIGO frequencies is only a bound on \(\eta\).

\paragraph{Black-hole images and ringdowns.}
Is there a visible deviation from Kerr at the photon ring or in
quasi-normal modes? That bounds \(r_0\). A perfectly Kerr object
is compatible with microscopic \(r_0\).

\paragraph{Galaxies, wide binaries, and clusters.}
Do disks accept a single \(\nu_{\sqrt{}}(g_{\mathrm{N}})\) with
\(\nu\to 1\) at Solar-System acceleration?
Appendix~D says yes, at RAR accuracy.
Do solar-neighbourhood wide binaries accept the isolated
evaluation of that law?
No: Assumption~\ref{ass:efe} saturates them at
\(\nu_{\sqrt{}}(g_{\mathrm{gal}})\).
Do merging clusters still require a stress that does not track
baryons?
That remains an open constraint on Assumption~\ref{ass:track}, not on
Theorem~7.3, and the external field is not a substitute for the
lensing offset.

\paragraph{Cosmology.}
Is \(\Lambda_{\mathrm{eff}}\) constant enough for a
Friedmann equation, and is the ultraviolet quartic absent from
the expansion history? This tests Theorem~\ref{thm:Leff} as a
mechanism, not as a number. A detected bulk \(\LUV^4\) vacuum
energy would kill the HBEV subtraction.

\paragraph{Tidal disruptions.}
Is there extra heating or altered fallback that hydrodynamics
plus GR tides cannot account for? A possible window on
Assumption~\ref{ass:drag}. A null result bounds \(\eta\) at
\(r_{\mathrm{t}}\) and nothing else.

\section{Claims the book has not made}
\label{sec:ch16-not-made}

The following sentences do not appear as theorems in this text,
and must not appear on a cover, in an abstract, or in a chapter
heading as if they did.

\begin{itemize}
  \item Dark matter has been explained.
  \item The cosmological-constant problem has been solved
        numerically.
  \item Singularities have been proved absent in dynamical
        collapse.
  \item LIGO inspirals are driven by a \(1/d^4\) drag.
  \item Lorentzian signature has been derived from a Euclidean
        matrix model.
  \item Presentation~A has been reduced to Presentation~B without
        Assumption~\ref{ass:dictionary}.
  \item The rationals \(A,B,C\) in the coefficient map are unique.
  \item Isolated \(\nu_{\sqrt{}}\) has been derived, and wide
    binaries have been predicted without an external field.
\end{itemize}

A sale-ready monograph that prints any of those sentences as
results has left this book and entered advertising.

\section{How to use Chapter~15 in public}
\label{sec:ch16-public}

Figure~\ref{fig:lab} may be shown as a numerical bookkeeping check of
an assumed drag in a Newtonian laboratory. It may not be
shown as ``NAMCG predicting an inspiral.'' The caption must
retain the words ``reduced model.'' The toy value
\(G_{\mathrm{toy}}=36\) must not be silently replaced by
\(G_{\mathrm{N}}\) in a talk slide.

\section{Exercises}
\label{sec:ch16-exercises}

\begin{enumerate}
  \item \textbf{Sort the list.}
    Take any popular article on emergent gravity and assign each
    of its claims to ``killer,'' ``coefficient,'' or ``not a
    claim of this book.''
  \item \textbf{One-page chain.}
    Write the eleven-line chain of Section~\ref{sec:ch16-chain}
    from memory, without looking. Where you stall is the section
    to reread.
  \item \textbf{Design a killer.}
    Propose one measurement, not already listed, that would
    refute Theorem~\ref{thm:mT} without passing through
    \(\eta\). State the null and the fatal outcome.
  \item \textbf{Design a non-killer.}
    Propose one measurement that looks dramatic and would only
    bound \(r_0\).
  \item \textbf{Abstract hygiene.}
    Rewrite a five-sentence abstract of this book that contains
    no sentence from Section~\ref{sec:ch16-not-made}.
\end{enumerate}

\section{Bridge to Chapter~17}
\label{sec:ch16-bridge}

The killers and the coefficients are now distinguished. What
remains is a research agenda: the calculations that would turn
Assumption~\ref{ass:dictionary} into a lemma, recompute
\(A,B,C\), evaluate \(S_{\mathrm{ent}}^{\mathrm{(IR)}}\), derive
\(\mu_{\mathrm{ent}}\) from \(I(A:B)\), including the external-field argument of Assumption~\ref{ass:efe}, and decide whether
Lorentzian signature can be more than Assumption~\ref{ass:lorentzian}.
Chapter~17 writes that agenda at the thickness of a reader who
has read the book and wants to work, not at the thickness of a
grant advertisement.